% Standalone problem document, generated from the adjacent README.md.
% Compile with XeLaTeX. Mathematical content is maintained in Markdown.
\documentclass[11pt,a4paper]{article}
\usepackage[fontsize=10.5pt]{scrextend}
\usepackage[margin=22mm,headheight=15pt,headsep=9mm,footskip=11mm]{geometry}
\usepackage{fontspec}
\setmainfont{texgyrepagella}[Extension=.otf,UprightFont=*-regular,BoldFont=*-bold,ItalicFont=*-italic,BoldItalicFont=*-bolditalic]
\setsansfont{texgyreheros}[Extension=.otf,UprightFont=*-regular,BoldFont=*-bold,ItalicFont=*-italic,BoldItalicFont=*-bolditalic]
\setmonofont{texgyrecursor}[Extension=.otf,UprightFont=*-regular,BoldFont=*-bold,ItalicFont=*-italic,BoldItalicFont=*-bolditalic,Scale=MatchLowercase]
\usepackage{amsmath,amssymb}
\usepackage{unicode-math}
\setmathfont{texgyrepagella-math.otf}
\usepackage{xcolor,microtype,enumitem,needspace,fancyhdr}
\usepackage{bookmark}
\definecolor{ink}{HTML}{17344B}
\definecolor{accent}{HTML}{197D80}
\definecolor{muted}{HTML}{596773}
\hypersetup{colorlinks=true,linkcolor=accent,urlcolor=accent,pdftitle={IE-23 - Uniqueness
of the right inverse minimizing an induced p-to-2
norm},pdfauthor={Open Problems in Numerical Linear Algebra}}
\urlstyle{same}
\setlength{\parindent}{0pt}
\setlength{\parskip}{5pt plus 1pt minus 1pt}
\linespread{1.04}
\setlength{\emergencystretch}{3em}
\widowpenalty=10000
\clubpenalty=10000
\displaywidowpenalty=10000
\allowdisplaybreaks[1]
\setlist{nosep,leftmargin=1.5em}
\providecommand{\tightlist}{\setlength{\itemsep}{0pt}\setlength{\parskip}{0pt}}
\providecommand{\pandocbounded}[1]{#1}
\pagestyle{fancy}
\fancyhf{}
\fancyhead[L]{\sffamily\footnotesize\color{muted}OPEN PROBLEMS IN NUMERICAL LINEAR ALGEBRA}
\fancyhead[R]{\sffamily\bfseries\footnotesize\color{accent}IE-23}
\fancyfoot[L]{\parbox[b]{0.9\textwidth}{\sffamily\fontsize{8}{10}\selectfont\color{muted}Editorial ratings; a bounded search cannot certify that a problem remains unsolved.\\Literature check: 2026-09-11}}
\fancyfoot[R]{\sffamily\footnotesize\color{muted}\thepage}
\renewcommand{\headrulewidth}{0.3pt}
\renewcommand{\headrule}{\hbox to\headwidth{\color{accent}\leaders\hrule height \headrulewidth\hfill}}
\makeatletter
\renewcommand{\section}{\Needspace{5\baselineskip}\@startsection{section}{1}{0pt}{12pt plus 2pt minus 2pt}{4pt}{\sffamily\bfseries\large\color{ink}}}
\renewcommand{\subsection}{\Needspace{4\baselineskip}\@startsection{subsection}{2}{0pt}{9pt plus 1pt minus 1pt}{3pt}{\sffamily\bfseries\normalsize\color{ink}}}
\makeatother
\setcounter{secnumdepth}{0}
\begin{document}
{\sffamily\small\color{muted}Linear systems and elimination\par}
\vspace{3pt}
{\raggedright\hyphenpenalty=10000\exhyphenpenalty=10000\sffamily\bfseries\fontsize{21}{24}\selectfont\color{ink}Uniqueness
of the right inverse minimizing an induced p-to-2 norm\par}
\vspace{7pt}
{\sffamily\small\textbf{Difficulty:} hard\quad\textbf{Importance:} interesting
to specialist\par}
{\sffamily\small\bfseries\color{accent}Status: Solved\par}
\vspace{4pt}
{\color{accent}\hrule height 0.8pt}
\vspace{6pt}
\subsection{Independently reviewed resolution -
2026-09-11}\label{independently-reviewed-resolution---2026-09-11}

\textbf{Negative resolution.} Theorem 1 gives a \(2\times3\)
full-row-rank matrix with distinct norm-minimizing right inverses for
every \(2<p<\infty\) over both fields. Their common induced norm is
\(2^{1/2-1/p}\). Sections 3-4 identify the complete minimizer sets,
including a complex higher-dimensional family. The smallest matrix is
attributed to Dokmanic and Gribonval\textquotesingle s spectral-norm
example; its extension to the displayed direct-inverse objective is
checked. No conclusion about the separate product objective is claimed.

\textbf{Author:} Matthew J. Colbrook, Department of Applied Mathematics
and Theoretical Physics, University of Cambridge.
\href{https://github.com/ajt60gaibb/OpenProblemsInNLA/blob/main/references/colbrook-recovered-2026-09-11/manuscripts/IE-23.pdf}{Complete
manuscript},
\href{https://github.com/ajt60gaibb/OpenProblemsInNLA/blob/main/references/colbrook-recovered-2026-09-11/verification/reviews/IE-23-review.md}{independent
proof review}, and
\href{https://github.com/ajt60gaibb/OpenProblemsInNLA/blob/main/references/colbrook-recovered-2026-09-11/README.md}{submission
record}. The supplied notes were reconstructed with substantial AI
assistance. This is independent agent verification, not external human
peer review or formal proof-assistant certification; no novelty or
priority claim is made.

The difficulty, importance and rating rationale below are historical
assessments of the original open target. Original statements, references
and dated audits are preserved.

\textbf{Rating rationale:} Hard reflects a focused uniqueness question
for a known norm minimizer; specialist impact is the characterization of
generalized inverses under induced norms.

\subsection{Problem statement}\label{problem-statement}

Let \(1\le m<n\), let \(A\in\mathbb C^{m\times n}\) have rank \(m\), and
let \(2<p<\infty\). For \(X\in\mathbb C^{n\times m}\) define the induced
norm

\[
\|X\|_{p\to2}=\sup_{y\in\mathbb C^m\setminus\{0\}}
\frac{\|Xy\|_2}{\bigl(\sum_{i=1}^m|y_i|^p\bigr)^{1/p}}.
\]

The Moore--Penrose inverse \(A^\dagger=A^*(AA^*)^{-1}\) minimizes this
norm among the right inverses \(X\) satisfying \(AX=I_m\). Is it always
the unique minimizer? Equivalently, must

\[
AX=I_m,\quad X\ne A^\dagger
\quad\Longrightarrow\quad
\|X\|_{p\to2}>\|A^\dagger\|_{p\to2}
\]

hold for every such \(A,p,X\)?

This is the direct-inverse uniqueness question in Dokmanić--Gribonval's
Remark 4.1. The endpoints and the separate objective \(\|XA\|_{p\to2}\)
are outside this statement.

\subsection{Connection to numerical linear
algebra}\label{connection-to-numerical-linear-algebra}

The norm measures the worst Euclidean solution amplification for
right-hand sides bounded in \(\ell^p\). Uniqueness determines whether
another linear solver can match the Moore--Penrose inverse's optimal
amplification while changing its sparsity or structure.

\newpage
\subsection{References}\label{references}

\begin{enumerate}
\def\labelenumi{\arabic{enumi}.}
\tightlist
\item
  I. Dokmanić and R. Gribonval, \emph{Beyond Moore--Penrose Part I:
  Generalized Inverses that Minimize Matrix Norms}, arXiv:1706.08349v2
  (2017), §2.1 and §2.4 (right inverses and \(A^\dagger\)); §4.4,
  Corollary 4.2(3) and Remark 4.1, manuscript p.~18 (PDF page 18).
  \href{https://arxiv.org/abs/1706.08349}{Preprint and version history}.
  \href{https://arxiv.org/pdf/1706.08349}{PDF}.
\item
  I. Dokmanić and R. Gribonval,
  \href{https://arxiv.org/abs/1706.08701}{\emph{Part II: The Sparse
  Pseudoinverse}}, 2017, §2. Its entrywise-norm objective is different.
\end{enumerate}

\subsection{Earlier status check ---
2026-09-08}\label{earlier-status-check-2026-09-08}

Checked the complete v2 statement and version history on 2026-09-08; v2,
dated 2017-07-13, remains the latest arXiv version of Part I. Searches
combined the authors, the paper title, ``induced norm'', ``p to 2'',
``unique'', ``uniqueness'', ``generalized inverse'', and 2025/2026. The
companion paper's sparse-inverse uniqueness results concern entrywise
norms and generic inputs, not this induced-norm claim for every
full-row-rank matrix. No later proof or counterexample was located. This
bounded search does not certify that no resolution exists elsewhere.

\subsection{Audit update --- 2026-09-10}\label{audit-update-2026-09-10}

The
\href{https://dokmanic.ece.illinois.edu/assets/pdf/DokmanicG17aa.pdf}{author
copy of Part I}, Corollary 4.2(3) and Remark 4.1 on printed p.~18,
explicitly separates minimality from the remaining uniqueness question
for \(2<p<\infty\). Searches for later induced-\(p\)-to-2 uniqueness
results found no resolution; Part II's sparse-inverse objectives do not
supply this missing assertion.
\end{document}
